\oldpage[III]{208}
\subsection{Application to proper morphisms: III. Invariance of the
Euler--Poincar\'e characteristic and the Hilbert polynomial}
\label{subsection:III.7.9}
\label{III.7.9}

\begin{env}[7.9.1]
\label{III.7.9.1}
Let $A$ be a ring and let $M$ be a projective $A$-module of
finite type.  Recall (Bourbaki, \emph{Alg.~comm.}, chap.~II,
\S5, no.~2) that this is equivalent to saying that the
$\sh{O}_X$-Module $\widetilde M$ associated with $M$ on
$X=\Spec(A)$ is locally free of finite type.  For every
$\mathfrak p\in\Spec(A)$, the \emph{rank of $M$ at
$\mathfrak p$}, denoted $\operatorname{rank}_{\mathfrak p}(M)$,
is the rank of the free $A_{\mathfrak p}$-module
$M_{\mathfrak p}$ (or, equivalently, the rank at
$\mathfrak p$ of the locally free $\sh{O}_X$-Module
$\widetilde M$).  Thus
\[
  \operatorname{rank}_{\mathfrak p}M
  =
  \operatorname{rank}_{\mathfrak p}(M_{\mathfrak p})
  =
  \operatorname{rank}_{\kres(\mathfrak p)}
    \bigl(M\otimes_A\kres(\mathfrak p)\bigr).
  \tag{7.9.1.1}\label{III.7.9.1.1}
\]
\end{env}

\begin{proposition}[7.9.2]
\label{III.7.9.2}
Let $P_\bullet$ be a finite complex of projective $A$-modules
of finite type, and, for every $A$-module $M$, set
\[
  T_\bullet(M)=H_\bullet(P_\bullet\otimes_A M).
\]
Then, for every $\mathfrak p\in\Spec(A)$,
\[
  \sum_i(-1)^i
    \operatorname{rank}_{\kres(\mathfrak p)}
      T_i\bigl(\kres(\mathfrak p)\bigr)
  =
  \sum_i(-1)^i\operatorname{rank}_{\mathfrak p}(P_i).
  \tag{7.9.2.1}\label{III.7.9.2.1}
\]
\end{proposition}

\begin{proof}
Indeed,
$T_i(\kres(\mathfrak p))
 =H_i(P_\bullet\otimes_A\kres(\mathfrak p))$ by definition.
In view of \eqref{III.7.9.1.1}, formula
\eqref{III.7.9.2.1} is simply the invariance of the
Euler--Poincar\'e characteristic of a finite complex of
finite-dimensional vector spaces upon passage to homology
\sref[0]{0.11.10.2}.
\end{proof}

\begin{corollary}[7.9.3]
\label{III.7.9.3}
The function
\[
  \mathfrak p\longmapsto
  \sum_i(-1)^i
    \operatorname{rank}_{\kres(\mathfrak p)}
      T_i\bigl(\kres(\mathfrak p)\bigr)
\]
is locally constant on $\Spec(A)$.
\end{corollary}

\begin{theorem}[7.9.4]
\label{III.7.9.4}
Let $Y$ be a locally Noetherian prescheme, let
$f:X\to Y$ be a proper morphism, and let
$\sh{P}_\bullet$ be a finite complex of coherent,
$Y$-flat $\sh{O}_X$-Modules.  If
\[
  \sh{T}_\bullet(\sh{M})
  =
  \mathscr{H}^{-\bullet}
    \bigl(f,\sh{P}^\bullet
      \otimes_{\sh{O}_Y}\sh{M}\bigr)
  \qquad\text{\rm(cf.~\eqref{III.7.7.1.1})},
\]
then the function
\oldpage[III]{209}
\[
  y\longmapsto
  \sum_i(-1)^i
    \operatorname{rank}_{\kres(y)}
      \sh{T}_i\bigl(\kres(y)\bigr)
  \tag{7.9.4.1}\label{III.7.9.4.1}
\]
is locally constant on $Y$.
\end{theorem}

\begin{proof}
We may restrict ourselves to the case where
$Y=\Spec(A)$ is affine and $A$ is Noetherian.  Since the
complex $\sh{P}_\bullet$ is finite, \sref{III.7.7.12}{(i)}
shows that
\[
  \sh{T}_p(\sh{M})
  =
  \mathscr{H}_p(\sh{L}_\bullet\otimes_{\sh{O}_Y}\sh{M}),
\]
where $\sh{L}_\bullet=\widetilde{L}_\bullet$ and
$L_\bullet$ is a finite complex of projective $A$-modules
of finite type.  The theorem now follows from
\sref{III.7.9.3}.
\end{proof}

\begin{env}[7.9.5]
\label{III.7.9.5}
Under the hypotheses of \sref{III.7.9.4}, the function
\eqref{III.7.9.4.1} is \emph{constant} when $Y$ is
\emph{connected}.  When $Y$ is connected and nonempty, its
unique (integral) value is denoted
\[
  \operatorname{EP}(f,\sh{P}_\bullet),\qquad
  \operatorname{EP}(Y,\sh{P}_\bullet),
  \qquad\text{or simply}\qquad
  \operatorname{EP}(\sh{P}_\bullet)
\]
when no confusion is possible; this integer is called the
\emph{Euler--Poincar\'e characteristic} of
$\sh{P}_\bullet$ relative to $f$ (or to $Y$).  In general,
we likewise write
\[
  \operatorname{EP}(f,\sh{P}_\bullet;y),\qquad
  \operatorname{EP}(Y,\sh{P}_\bullet;y),\qquad
  \operatorname{EP}(\sh{P}_\bullet;y)
\]
for the right-hand side of \eqref{III.7.9.4.1}.
\end{env}

\begin{env}[7.9.6]
\label{III.7.9.6}
Under the hypotheses of \sref{III.7.9.4} concerning
$X$, $Y$, and $f$, let
\[
  0\longrightarrow\sh{P}'_\bullet
    \xrightarrow{u}\sh{P}_\bullet
    \xrightarrow{v}\sh{P}''_\bullet
    \longrightarrow 0
\]
be an exact sequence of finite complexes of coherent,
$Y$-flat $\sh{O}_X$-Modules, where the homomorphisms $u$
and $v$ have even degrees $2d$ and $2d'$, respectively.
Since $\sh{T}_\bullet$ is a homological functor
\sref{III.7.7.1}, there is an exact homology sequence
\[
  \cdots\longrightarrow
  \sh{T}_i(\sh{P}'_\bullet,\kres(y))
  \longrightarrow
  \sh{T}_{i+2d}(\sh{P}_\bullet,\kres(y))
  \longrightarrow
  \sh{T}_{i+2d+2d'}(\sh{P}''_\bullet,\kres(y))
  \longrightarrow
  \sh{T}_{i-1}(\sh{P}'_\bullet,\kres(y))
  \longrightarrow\cdots ,
\]
which moreover has only finitely many terms.  On writing
that the Euler--Poincar\'e characteristic of this complex is
zero \sref[0]{0.11.10.1}, we obtain at once
\[
  \operatorname{EP}(\sh{P}_\bullet;y)
  =
  \operatorname{EP}(\sh{P}'_\bullet;y)
  +
  \operatorname{EP}(\sh{P}''_\bullet;y)
  \tag{7.9.6.1}\label{III.7.9.6.1}
\]
for every $y\in Y$.

Suppose, for example, that
$\sh{P}_\bullet=(\sh{P}_i)$ and that $\sh{P}_i=0$ for
$i<0$.  We then have the following exact sequence of
complexes:
\[
\xymatrix@C=2.4em{
  \cdots \ar[r] & 0 \ar[r]\ar[d] & 0 \ar[r]\ar[d]
    & 0 \ar[r]\ar[d] & 0 \ar[r] & \cdots \\
  \cdots \ar[r] & 0 \ar[r]\ar[d] & 0 \ar[r]\ar[d]
    & \sh{P}_1 \ar[r]\ar[d] & \sh{P}_2 \ar[r] & \cdots \\
  \cdots \ar[r] & 0 \ar[r]\ar[d] & \sh{P}_0 \ar[r]\ar[d]
    & \sh{P}_1 \ar[r]\ar[d] & \sh{P}_2 \ar[r] & \cdots \\
  \cdots \ar[r] & 0 \ar[r]\ar[d] & \sh{P}_0 \ar[r]\ar[d]
    & 0 \ar[r]\ar[d] & 0 \ar[r] & \cdots \\
  \cdots \ar[r] & 0 \ar[r] & 0 \ar[r]
    & 0 \ar[r] & 0 \ar[r] & \cdots
}
\]
the nonzero vertical arrows being identity automorphisms.
Applying \eqref{III.7.9.6.1} to this exact sequence and
arguing by induction on the length of $\sh{P}_\bullet$
gives
\[
  \operatorname{EP}(\sh{P}_\bullet;y)
  =
  \sum_i(-1)^i\operatorname{EP}(\sh{P}_i;y).
  \tag{7.9.6.2}\label{III.7.9.6.2}
\]

\oldpage[III]{210}
Here, for each coherent $\sh{O}_X$-Module $\sh{F}$ that is
flat over $Y$, the notation
$\operatorname{EP}(\sh{F};y)$ (or
$\operatorname{EP}(f,\sh{F};y)$, or
$\operatorname{EP}(Y,\sh{F};y)$) means
$\operatorname{EP}(\sh{Q}_\bullet;y)$ (with the
corresponding relative notation) for the complex
$\sh{Q}_\bullet$ whose only nonzero term is $\sh{F}$ in
degree $0$.  Thus the study of Euler--Poincar\'e
characteristics of complexes reduces to that of complexes
with a single term.
\end{env}

\begin{proposition}[7.9.7]
\label{III.7.9.7}
Under the hypotheses of \sref{III.7.9.4}, let $Y'$ be a
locally Noetherian prescheme, let $g:Y'\to Y$ be a
morphism, put
\[
  X'=X\times_Y Y',\qquad
  f'=f_{(Y')}:X'\longrightarrow Y',
\]
and let
\[
  \sh{P}'_\bullet
  =
  \sh{P}_\bullet\otimes_{\sh{O}_Y}\sh{O}_{Y'}
\]
be the resulting finite complex of $\sh{O}_{X'}$-Modules.
Then $\sh{P}'_\bullet$ consists of coherent,
$Y'$-flat $\sh{O}_{X'}$-Modules, and, for every
$y'\in Y'$,
\[
  \operatorname{EP}(\sh{P}'_\bullet;y')
  =
  \operatorname{EP}(\sh{P}_\bullet;g(y')).
  \tag{7.9.7.1}\label{III.7.9.7.1}
\]
\end{proposition}

\begin{proof}
The $\sh{O}_{X'}$-Modules $\sh{P}'_i$, being inverse images
of the $\sh{P}_i$ under the projection $X'\to X$, are
coherent.  They are $Y'$-flat by \sref[0]{0.6.2.1} and
\sref[I]{I.4.14.5}, the question being local on $X$, $Y$,
and $Y'$.  Finally, $f'$ is proper by
\sref[II]{II.5.4.2}, so the left-hand side of
\eqref{III.7.9.7.1} is defined.  Formula
\eqref{III.7.9.7.1} now follows from
\textup{(6.10.4.2)}, \sref{III.7.7.2}, and
\sref{III.7.6.7}, after reducing, as one may always do, to
the case where $Y$ and $Y'$ are affine.
\end{proof}

\begin{proposition}[7.9.8]
\label{III.7.9.8}
Suppose that the hypotheses of \sref{III.7.9.4} hold and,
in addition, that there is an integer $i_0$ such that
\[
  \sh{T}_i(\kres(y))=0
  \qquad\text{for $i\neq i_0$ and every $y\in Y$.}
\]
Then
\[
  \sh{T}_{i_0}(\sh{O}_Y)
  =
  \mathscr{H}^{-i_0}(f,\sh{P}^\bullet)
\]
is a locally free $\sh{O}_Y$-Module, and its rank at
$y\in Y$ is
\[
  (-1)^{i_0}\operatorname{EP}(f,\sh{P}_\bullet;y).
\]
\end{proposition}

\begin{proof}
First observe that the hypotheses of \sref{III.7.4.4} are
satisfied by the functors $\sh{T}^{\sh{O}_y}_i$, so that
\sref{III.7.4.7} applies to them.  The hypothesis and
\sref{III.7.5.3} imply that
$\sh{T}^{\sh{O}_y}_i$ is zero for $i\neq i_0$.
Consequently $\sh{T}_{i_0}$ is exact by
\sref{III.7.3.3}, and \sref{III.7.8.4} shows that
$\mathscr{H}^{-i_0}(f,\sh{P}^\bullet)$ is locally free.
Its rank at $y\in Y$ is
\[
  \operatorname{rank}_{\kres(y)}
    \sh{T}_{i_0}(\kres(y))
  =
  (-1)^{i_0}\operatorname{EP}(f,\sh{P}_\bullet;y),
\]
by definition, since $\sh{T}_i(\kres(y))=0$ for
$i\neq i_0$.\footnote{The printed French proof omits the
factor $(-1)^{i_0}$ in this displayed equality; it is
required both by the definition and by the statement of
the proposition.}
\end{proof}

\begin{corollary}[7.9.9]
\label{III.7.9.9}
Let $Y$ be a locally Noetherian prescheme, let
$f:X\to Y$ be a proper morphism, and let $\sh{F}$ be a
coherent, $Y$-flat $\sh{O}_X$-Module.  Suppose that there
is an integer $i_0$ such that
\[
  H^i\bigl(f^{-1}(y),
    \sh{F}\otimes_{\sh{O}_Y}\kres(y)\bigr)=0
\]
for every $i\neq i_0$ and every $y\in Y$.  Then
$R^{i_0}f_*(\sh{F})$ is a locally free
$\sh{O}_Y$-Module, and its rank at $y$ is
$(-1)^{i_0}\operatorname{EP}(f,\sh{F};y)$.
\end{corollary}

In particular:

\begin{corollary}[7.9.10]
\label{III.7.9.10}
Under the preliminary hypotheses of \sref{III.7.9.9}
concerning $X$, $Y$, and $\sh{F}$, suppose that
$R^if_*(\sh{F})=0$ for every $i>0$.  Then
$f_*(\sh{F})$ is a locally free $\sh{O}_Y$-Module, and its
rank at $y$ is $\operatorname{EP}(f,\sh{F};y)$.
\end{corollary}

\begin{proof}
In view of \sref{III.7.9.9}, it is enough to prove the
following lemma.
\end{proof}

\begin{lemma}[7.9.10.1]
\label{III.7.9.10.1}
Under the hypotheses of \sref{III.7.9.10},
\[
  H^i\bigl(f^{-1}(y),
    \sh{F}\otimes_{\sh{O}_Y}\kres(y)\bigr)=0
\]
for every $i>0$ and every $y\in Y$.
\end{lemma}

\begin{proof}
We may restrict ourselves to the case where
$Y=\Spec(A)$ is affine.  With the notation of
\sref{III.7.9.4}, and with $\sh{P}_\bullet$ reduced to its
only nonzero term $\sh{F}$ in degree $0$, the hypothesis
gives $\sh{T}_p(\sh{O}_Y)=0$ for $p<0$.  It follows from
\sref{III.7.3.7} that $\sh{T}_p$ is exact for $p<0$, and
the lemma then follows from the equivalence of conditions
a) and d) in \sref{III.7.7.5}.
\end{proof}

\oldpage[III]{211}
\begin{proposition}[7.9.11]
\label{III.7.9.11}
Under the hypotheses of \sref{III.7.9.4}, let $\sh{L}$ be
an invertible $\sh{O}_X$-Module that is very ample over
$Y$, and, for every $n\in\bb{Z}$, put
\[
  \sh{P}_\bullet(n)
  =
  \sh{P}_\bullet
    \otimes_{\sh{O}_X}\sh{L}^{\otimes n}.
\]
Then, for every $y\in Y$, the function
\[
  n\longmapsto
  \operatorname{EP}(f,\sh{P}_\bullet(n);y)
  \tag{7.9.11.1}\label{III.7.9.11.1}
\]
is a polynomial with rational coefficients, and is the
same for all points in any one connected component of
$Y$.
\end{proposition}

\begin{proof}
It is clear that $\sh{P}_\bullet(n)$ is a complex of
$Y$-flat $\sh{O}_X$-Modules.  By
\eqref{III.7.9.6.2}, we may reduce to the case where
$\sh{P}_\bullet$ has a single nonzero term
$\sh{F}\neq0$ in degree $0$.  Since the questions are
local on $Y$, we may moreover suppose that $Y$ is affine
and that $f$ is projective \sref[II]{II.5.5.3}.  Put
\[
  X_y=f^{-1}(y),\qquad
  \sh{L}_y=\sh{L}\otimes_{\sh{O}_Y}\kres(y);
\]
then $\sh{L}_y$ is a very ample $\sh{O}_{X_y}$-Module
\sref[II]{II.4.4.10}.  By \sref{III.7.7.2}, for the
functor $\sh{T}_\bullet$ associated with
$\sh{P}_\bullet(n)$ we have
\[
  \sh{T}_i(\kres(y))
  =
  H^{-i}\bigl(
    X_y,\sh{F}_y\otimes\sh{L}_y^{\otimes n}\bigr),
  \qquad
  \sh{F}_y=\sh{F}\otimes_{\sh{O}_Y}\kres(y).
\]
Thus $\operatorname{EP}(f,\sh{F}(n);y)$ is precisely the
Euler--Poincar\'e characteristic
$\chi_{\kres(y)}(\sh{F}_y(n))$ defined in
\sref{III.2.5.1}.  It follows from \sref{III.2.5.3} that
\eqref{III.7.9.11.1} is a polynomial.  For each fixed
$n$, its value is constant on each connected component of
$Y$ by \sref{III.7.9.4}; this completes the proof.
\end{proof}

\begin{env}[7.9.12]
\label{III.7.9.12}
The polynomial \eqref{III.7.9.11.1}, with rational
coefficients, will be denoted
\[
  \operatorname{PH}(f,\sh{P}_\bullet;y)
  \qquad\text{or}\qquad
  \operatorname{PH}(\sh{P}_\bullet;y),
\]
and will be called the \emph{Hilbert polynomial at $y$}
relative to $\sh{P}_\bullet$, $f$, and $\sh{L}$ (or simply
the \emph{Hilbert polynomial at $y$} of
$\sh{P}_\bullet$, or of $f$, when no confusion is
possible).  When $Y$ is connected and nonempty, the
reference to $y$ is omitted from the notation and
terminology.  This invariant will play an essential role
in Chapter~V, in the theory of ``moduli'' of coherent
quotient sheaves of a given coherent sheaf.

With the notation of \sref{III.7.9.6} and
\sref{III.7.9.11}, we have
\[
  \operatorname{PH}(\sh{P}_\bullet;y)
  =
  \operatorname{PH}(\sh{P}'_\bullet;y)
  +
  \operatorname{PH}(\sh{P}''_\bullet;y),
  \tag{7.9.12.1}\label{III.7.9.12.1}
\]
and in particular
\[
  \operatorname{PH}(\sh{P}_\bullet;y)
  =
  \sum_i(-1)^i\operatorname{PH}(\sh{P}_i;y).
  \tag{7.9.12.2}\label{III.7.9.12.2}
\]
These formulas follow immediately from
\eqref{III.7.9.6.1} and \eqref{III.7.9.6.2}.  Likewise,
with the notation and hypotheses of \sref{III.7.9.7},
\[
  \operatorname{PH}(\sh{P}'_\bullet;y')
  =
  \operatorname{PH}(\sh{P}_\bullet;g(y')).
  \tag{7.9.12.3}\label{III.7.9.12.3}
\]
\end{env}

Formula \eqref{III.7.9.12.2} reduces the study of the
Hilbert polynomials of a complex to that of the Hilbert
polynomial of a single $Y$-flat $\sh{O}_X$-Module.  Such
polynomials have the following remarkable interpretation,
independent of homological considerations.

\begin{corollary}[7.9.13]
\label{III.7.9.13}
Let $Y$ be a Noetherian prescheme, let $f:X\to Y$ be a
proper morphism, let $\sh{L}$ be an invertible
$\sh{O}_X$-Module that is very ample over $Y$, and let
$\sh{F}$ be a coherent, $Y$-flat $\sh{O}_X$-Module.  There
is an integer $n_0$ such that, for $n\geq n_0$,
$f_*(\sh{F}(n))$ is a locally free $\sh{O}_Y$-Module whose
rank at $y\in Y$ is
\[
  \operatorname{PH}(f,\sh{F};y)(n).
\]
\end{corollary}

\begin{proof}
\oldpage[III]{212}
Since $f$ is projective \sref[II]{II.5.5.3}, there is an
integer $n_0$ such that, for $n\geq n_0$,
\[
  R^if_*(\sh{F}(n))=0
  \qquad\text{for every $i>0$}
\]
by \sref{III.2.2.1}.  The conclusion follows from
\sref{III.7.9.10}.
\end{proof}

The following flatness criterion will be important in
Chapter~V, in the theory of ``moduli'' of coherent
sheaves.

\begin{proposition}[7.9.14]
\label{III.7.9.14}
Let $Y$ be a locally Noetherian prescheme, let
$f:X\to Y$ be a projective morphism, and let $\sh{L}$ be
an invertible $\sh{O}_X$-Module that is ample for $f$.
For every $\sh{O}_X$-Module $\sh{F}$ and every
$n\in\bb{Z}$, put
\[
  \sh{F}(n)=\sh{F}\otimes\sh{L}^{\otimes n}.
\]
A coherent $\sh{O}_X$-Module $\sh{F}$ is $Y$-flat if and
only if there is an integer $n_0$ such that, for every
$n\geq n_0$, $f_*(\sh{F}(n))$ is a locally free
$\sh{O}_Y$-Module.
\end{proposition}

\begin{proof}
The necessity of the condition is proved as in
\sref{III.7.9.13}; the result of \sref{III.2.2.1} applies
to an ample sheaf because $f$ is projective.

Conversely, we may reduce to the case where
$Y=\Spec(A)$ is affine.  By the hypothesis and
\sref{III.2.2.2}{(i)}, the $A$-modules
$\Gamma(X,\sh{F}(n))$ are projective and of finite type
(Bourbaki, \emph{Alg.~comm.}, chap.~II, \S5, no.~2,
th.~1).  Let
\[
  S=\bigoplus_{n\geq0}\Gamma(X,\sh{L}^{\otimes n})
\]
be the associated graded ring.  We know that $X$ is
canonically identified with $\Proj(S)$ by
\sref[II]{II.4.5.2}{(b)} and \sref{III.5.4.4}.  Let
\[
  M=\bigoplus_{n\geq n_0}\Gamma(X,\sh{F}(n)).
\]
After replacing $\sh{L}$, if necessary, by a power
$\sh{L}^{\otimes d}$, we may suppose that $S$ is
generated by finitely many elements of degree $1$
\sref{III.2.3.5.1}; it then follows from
\sref[II]{II.2.7.5} and \sref[II]{II.2.7.2} that $\sh{F}$ is
identified with $\shProj_0(M)$.

For every homogeneous element $g\in S$ of positive
degree, we therefore have
\[
  \Gamma(X_g,\sh{F})=M_{(g)}.
\]
Now $M$, being a direct sum of projective $A$-modules, is
a flat $A$-module.  Hence $M_g$ is flat by
\sref[0]{0.6.3.2}, and so is $M_{(g)}$, which is a direct
factor of $M_g$ by \sref[0]{0.6.1.2}.  We conclude from
\sref[I]{I.4.14.5} that $\sh{F}$ is $Y$-flat at every
point of $X_g$.  Since the open subsets $X_g$ cover $X$,
the proposition is proved.
\end{proof}

\bigskip
\noindent\emph{(To be continued.)}

\clearpage
\oldpage[III]{213}
\section*{INDEX OF NOTATION}

\begin{flushleft}
$\mathbf{Tor}_n^A(P_\bullet,Q_\bullet)$: 6.3.1.

\medskip
$\shTor_n^{\sh{O}_S}(\sh{P}_\bullet,\sh{Q}_\bullet)$,
$\shTor_n^S(\sh{P}_\bullet,\sh{Q}_\bullet)$: 6.4.1 and 6.5.1.

\medskip
$\shTor_n^S(\sh{F},\sh{G})$: 6.5.1.

\medskip
$\shTor_q^S(\sh{P}_\bullet^{(1)},\sh{P}_\bullet^{(2)},\ldots,
\sh{P}_\bullet^{(m)})$: 6.5.15.

\medskip
$\mathbf H^{-n}(\mathfrak U^{(i)},\sh{P}_\bullet^{(i)})$,
$\mathbf H^{-n}(X^{(i)},\sh{P}_\bullet^{(i)})$: 6.6.1.

\medskip
$\mathbf{Tor}_n^A(L_{\bullet\bullet}^{(1)},L_{\bullet\bullet}^{(2)})$,
$\mathbf{Tor}_n^S(\mathfrak U^{(1)},\mathfrak U^{(2)};
\sh{P}_\bullet^{(1)},\sh{P}_\bullet^{(2)})$: 6.6.2.

\medskip
$\shTor_n^S(\mathfrak U^{(1)},\mathfrak U^{(2)};
\sh{F}^{(1)},\sh{F}^{(2)})$,
$\shTor_n^S(\mathfrak U^{(1)},\mathfrak U^{(2)};
\sh{P}_\bullet^{(1)},\sh{P}_\bullet^{(2)})$: 6.6.2.

\medskip
$\shTor_n^S(f_1,f_2;\sh{P}_\bullet^{(1)},
\sh{P}_\bullet^{(2)})$: 6.6.8, 6.7.1, and 6.7.3.

\medskip
$\shTor_n^S((f_i)_{i\in I};(\sh{P}_\bullet^{(i)})_{i\in I})$,
$\shTor_n^S(f_1,\ldots,f_m;\sh{P}_\bullet^{(1)},\ldots,
\sh{P}_\bullet^{(m)})$: 6.7.11.

\medskip
$\mathbf{Ab}$, $\mathbf{Ab}_A$: 7.1.1.

\medskip
$T_{(B)}$, $T^{(B)}$, $T\otimes_A B$
($T$ a functor from $\mathbf{Ab}_A$ to $\mathbf{Ab}$): 7.1.3.

\medskip
$t_M$: 7.2.2.

\medskip
$T_{\mathfrak p}$: 7.1.4.

\medskip
$\sh{T}_\bullet^{Y'}(\sh{M}')$,
$\sh{T}_\bullet^U(\sh{M}|U)$: 7.7.2.

\medskip
$\operatorname{EP}(f,\sh{P}_\bullet;y)$,
$\operatorname{EP}(Y,\sh{P}_\bullet;y)$,
$\operatorname{EP}(\sh{P}_\bullet;y)$: 7.9.5.

\medskip
$\operatorname{PH}(f,\sh{P}_\bullet;y)$,
$\operatorname{PH}(\sh{P}_\bullet;y)$: 7.9.12.
\end{flushleft}

\clearpage
\oldpage[III]{214}
\section*{TERMINOLOGICAL INDEX}

\begin{flushleft}
Euler--Poincaré characteristic of a complex of Modules: 7.9.5.

\medskip
Cohomologically flat at a point $y\in Y$ in dimension $p$,
cohomologically flat over $Y$ in dimension $p$, cohomologically flat over
$Y$ at the point $y$, cohomologically flat over $Y$ (complex of
$\sh{O}_X$-Modules flat over $Y$): 7.8.1.

\medskip
Homotopic complexes: 6.1.4.

\medskip
Ringed space of cohomological dimension $\leq n$: 6.5.5.

\medskip
Extension of scalars in an additive covariant functor: 7.1.3.

\medskip
Sheaf of rings of cohomological dimension $\leq n$: 6.5.5.

\medskip
Künneth formula: 6.7.8.

\medskip
Homologically flat at a point $y\in Y$ in dimension $p$, homologically flat
over $Y$ in dimension $p$, homologically flat over $Y$ at the point $y$,
homologically flat over $Y$ (complex of $\sh{O}_X$-Modules flat over $Y$):
7.8.1.

\medskip
Homotopy: 6.1.4.

\medskip
Hypertor of two complexes of $A$-modules: 6.3.1.

\medskip
Local hypertor of two complexes of Modules: 6.4.1.

\medskip
Global hypertor of two complexes of Modules: 6.6.2, 6.6.8, and 6.7.3.

\medskip
Hilbert polynomial relative to a complex of Modules: 7.9.12.

\medskip
Künneth spectral sequences: 6.7.3.
\end{flushleft}

\clearpage
\oldpage[III]{215}
\section*{TABLE OF CONTENTS}

\begin{flushleft}
\textsc{Chapter III.} --- \textbf{Cohomological study of coherent sheaves
(continued)} \dotfill 5

\medskip
\S~6. Local and global Tor functors; Künneth formula \dotfill 5

6.1. Introduction \dotfill 5

6.2. Hypercohomology of complexes of Modules on a prescheme \dotfill 6

6.3. Hypertor of two complexes of modules \dotfill 9

6.4. Local hypertor functors of complexes of quasi-coherent Modules; the case
of affine schemes \dotfill 14

6.5. Local hypertor functors of complexes of quasi-coherent Modules: the
general case \dotfill 16

6.6. Global hypertor functors of complexes of quasi-coherent Modules and
Künneth spectral sequences: the case of an affine base \dotfill 21

6.7. Global hypertor functors of complexes of quasi-coherent Modules and
Künneth spectral sequences: the general case \dotfill 25

6.8. Associativity spectral sequences for global hypertor \dotfill 32

6.9. Base-change spectral sequences in global hypertor \dotfill 34

6.10. Local structure of certain cohomological functors \dotfill 39

\medskip
\S~7. Study of base change in covariant homological functors of Modules
\dotfill 43

7.1. Functors of $A$-modules \dotfill 43

7.2. Characterization of the tensor-product functor \dotfill 44

7.3. Exactness criteria for homological functors of modules \dotfill 48

7.4. Exactness criteria for the functors $\mathrm H_i(P\otimes_A M)$
\dotfill 53

7.5. The case of Noetherian local rings \dotfill 58

7.6. Descent of exactness properties. Semicontinuity theorem and Grauert's
exactness criterion \dotfill 60

7.7. Application to proper morphisms: I. The exchange property \dotfill 65

7.8. Application to proper morphisms: II. Cohomological flatness criteria
\dotfill 72

7.9. Application to proper morphisms: III. Invariance of the Euler--Poincaré
characteristic and the Hilbert polynomial \dotfill 76

\begin{flushright}
\emph{Received December 15, 1962.}
\end{flushright}
\end{flushleft}

% Printed page 216 is blank.
\clearpage
\thispagestyle{empty}
\null
\clearpage

\oldpage[III]{217}
\section*{ERRATA AND ADDENDA}
\begin{center}
(List 2)

\medskip
A) \emph{Typographical errors}
\end{center}

\noindent\textbf{$(0_{\mathrm I},\,3.2.1)$} On line 2 of p. 27, replace $X$
by $U$.

\noindent\textbf{$(0_{\mathrm I},\,3.2.6)$} On line 1 from the bottom of p.
27 and line 1 of p. 28, replace (3 times) $\mathfrak S_\lambda$ by
$\sh{H}_\lambda$.

\noindent\textbf{$(0_{\mathrm I},\,3.4.5)$} Add a parenthesis before 3.4.5).

\noindent\textbf{$(0_{\mathrm I},\,4.2.1)$} On line 10 of p. 39, replace
$\psi_*(A)$ by $\psi_*(\sh{A})$.

\noindent\textbf{$(0_{\mathrm I},\,5.1.3)$} On line 16 of p. 45, replace
$\sh{F}|V$ by $\sh{F}|U$.

\noindent\textbf{$(0_{\mathrm I},\,5.3.9)$} On line 12 of p. 48, before
``such that'', add: over an open neighborhood $U$ of $x$.

\noindent\textbf{$(0_{\mathrm I},\,7.6.9)$} On line 3 from the bottom of p.
73, replace $I_\lambda$ by $J_\lambda$.

\noindent\textbf{(I, 1.1.15)} On line 3 of p. 83, replace $\mathfrak a$ by
$\mathfrak a\ne A$.

\noindent\textbf{(I, 1.4.1)} On line 2 from the bottom of p. 90, replace
$\widetilde N$ by $N$.

\noindent\textbf{(I, 2.3.1)} On line 3 from the bottom of p. 101, replace (0,
4.1.6) by (0, 4.1.7).

\noindent\textbf{(I, 2.3.2)} On line 11 of p. 101, replace $B$ by $B_s$ and
$C$ by $C_t$.

\noindent\textbf{(I, 2.5.5)} On line 19 of p. 104, replace $S$-morphism by
$S$-prescheme.

\noindent\textbf{(I, 3.3.7)} On line 17 of p. 109, replace $Y_{(S)}$ by
$Y_{(S')}$.

\noindent\textbf{(I, 3.4.5)} On line 9 of p. 113, replace $\mathrm s$ by $s$.

\noindent\textbf{(I, 3.4.8)} On line 4 of p. 114, replace $p(x)$ by $f(x)$.

\noindent\textbf{(I, 3.5.1)} On line 3 of p. 115, replace $1_X\times g$ by
$1_{X'}\times g$.

\noindent\textbf{(I, 3.7.2)} On line 8 of p. 119, replace the first $X$ by
$X'$.

\noindent\textbf{(I, 4.2.2)} On line 8 of p. 123, in the left vertical arrow
of the diagram, replace $\alpha_{\psi(y)}$ by $\vphi_{\psi(y)}$.

\noindent\textbf{(I, 4.4.3)} On line 16 of p. 126, replace \emph{isomorphée}
by \emph{isomorphes}.

\noindent\textbf{(I, 4.5.5)} On line 17 from the bottom of p. 127, replace
(4.2.4) by (4.2.5).

\noindent\textbf{(I, 5.1.4)} On line 14 from the bottom of p. 128, replace
(2.1.7) by (2.1.8).

\noindent\textbf{(I, 5.3.13)} On line 20 of p. 134, replace (4.2.4) by
(4.2.5).

\noindent\textbf{(I, 6.4.2)} On line 1 from the bottom of p. 147, replace
\emph{covering} by \emph{finite covering}.

\noindent\textbf{(I, 9.1.13)} On line 15 of p. 171, replace
$p^{-1}(\sh{F})$ by $p^*(\sh{F})$.

\noindent\textbf{(I, 9.5.11)} On line 7 of p. 179, replace $Y$ by $Y'$.

\noindent\textbf{(I, 9.6.5)} On line 17 from the bottom of p. 180, replace
$(0,\,4.1.4)$ by $(0,\,4.1.3)$.

\noindent\textbf{(I, 10.12.2)} On line 14 from the bottom of p. 206, replace
$(X_n)_{(S_n)}$ by $(X_n)_{(S_m)}$.

\clearpage
\oldpage[III]{218}
\noindent\textbf{(I), Terminological index:} On line 6 from the bottom of p.
219, replace 0, 4.1.4 by 0, 4.1.5. On lines 16, 17, and 18 from the bottom of
p. 222, replace I, 2.1.7 by I, 2.1.8; on line 19 from the bottom of p. 222,
replace I, 2.1.8 by I, 2.1.7.

\noindent\textbf{(II, 1.5.2)} On line 17 of p. 12, insert an $f$ beside the
left vertical arrow of the diagram.

\noindent\textbf{(II, 4.2.7)} On line 16 from the bottom of p. 75, replace
(III, 2.1.14) by (III, 2.1.13).

\noindent\textbf{(II, 5.2.1)} On line 16 from the bottom of p. 97, replace
$X-U$ by $U$. On line 8 from the bottom of p. 97, replace $X'$ by $X_f$.

\noindent\textbf{(II, 5.3.6)} On line 16 of p. 100, replace (4.6.18) by
(4.6.17).

\noindent\textbf{(II, 6.2.7)} On line 9 from the bottom of p. 116, replace
chap. V by chap. IV.

\noindent\textbf{(II, 6.6.5)} On line 15 from the bottom of p. 132, replace
chap. V by chap. IV.

\noindent\textbf{(II, 7.4.12)} On line 1 of p. 151, replace chap. V by chap.
III.

\noindent\textbf{(II, 8.1.4)} On line 15 of p. 154, replace (III, 2.3.8) by
(III, 2.3.7).

\noindent\textbf{(II, 8.10.5)} On line 17 from the bottom of p. 187, replace
$g_{(j)x}$ by $g_{j(x)}$.

\noindent\textbf{$(0_{\mathrm{III}},\,10.2.7)$} On line 1 of p. 20, replace
$u$-adic by $\mathfrak m$-adic. On line 13 from the bottom of p. 20, replace
$k$ by $k_i$.

\noindent\textbf{$(0_{\mathrm{III}},\,11.1.1)$} On line 18 of p. 23, replace
$\operatorname{Im}(d_r^{p+r,q-r+1})$ by
$\operatorname{Im}(d_r^{p-r,q+r-1})$.

\noindent\textbf{$(0_{\mathrm{III}},\,11.8.6)$} On line 14 from the bottom of
p. 45, replace the arrows $\to$ by $\rightsquigarrow$ (4 times).

\noindent\textbf{$(0_{\mathrm{III}},\,12.3.4)$} On line 4 of p. 61, replace
$\rightsquigarrow$ by $\to$.

\noindent\textbf{$(0_{\mathrm{III}},\,13.2.4)$} On line 2 from the bottom of
p. 67, replace $\varinjlim$ by $\varprojlim$ (2 times).

\noindent\textbf{(III, 1.1.7)} On line 15 of p. 85, replace $A^r$ by
$\bigwedge(A^r)$ (2 times).

\noindent\textbf{(III, 1.2.4)} On line 7 of p. 87, replace
$\Gamma(U,\sh{F})$ by $\Gamma(U,\sh{O}_X)$ and replace $\sh{G}$ by $\sh{F}$.
On line 8 of p. 87, replace $\HH^p(\mathfrak U,\sh{G})$ by
$\HH^p(\mathfrak U,\sh{F})$.

\noindent\textbf{(III, 2.5.3.1)} On line 3 of p. 111, replace $n>0$ by
$n\geq0$. On line 8 of p. 111, replace $-r\leq n\leq0$ by $-r\leq n<0$.

\noindent\textbf{(III, 3.1.2)} On line 13 of p. 116, replace $\sh{G}$ by
$\sh{G}\in\mathbf K'$.

\noindent\textbf{(III, 4.3.6)} On line 7 of p. 133, replace
$K'=K\otimes_A\widehat A$ by $K'\supset K\otimes_A\widehat A$. On line 16 of
p. 133, replace $R(X)$ by $R(Y)$.

\noindent\textbf{(III, 4.4.12)} On line 10 of p. 138, replace chap. V by
chap. IV.

\noindent\textbf{(III, 4.6.6)} On line 19 of p. 142, replace ``chap. V'' by
``in a later section''.

\begin{center}
B) \emph{Textual modifications}\footnote{To facilitate the search for
references, modifications belonging to the errata list inserted in chapter N
will henceforth be designated by $\mathbf{Err}_{\mathrm N}$ followed by a
number.}
\end{center}

\noindent\textbf{$(\mathbf{Err}_{\mathrm{III}},\,1)$} In
$(0_{\mathrm I},\,5.1.3)$, on line 18 of p. 45, replace ``direct sum'' by
``finite direct sum''.

\noindent\textbf{$(\mathbf{Err}_{\mathrm{III}},\,2)$} In
$(0_{\mathrm I},\,5.4.3)$, on lines 8 to 4 from the bottom of p. 49, replace
the text beginning ``In the general case\ldots{}'' by the following:

In the general case where $X$ is a ringed space such that $\sh{O}_x$ is a
\emph{local} ring for
\oldpage[III]{219}
every $x\in X$, if $\sh{L}$ is an $\sh{O}_X$-Module \emph{of finite type} for
which there exists an $\sh{O}_X$-Module $\sh{F}$ such that
$\sh{L}\otimes_{\sh{O}_X}\sh{F}$ is \emph{isomorphic to} $\sh{O}_X$, the
preceding argument shows that, for every $x\in X$, $\sh{L}_x$ is isomorphic to
$\sh{O}_x$. It follows that $\sh{L}$ is then \emph{invertible}. Indeed, for
every $x\in X$, let $U$ be an open neighborhood of $x$ such that $\sh{L}|U$ is
generated by $n$ sections $s_i\ (1\leq i\leq n)$ over $U$ (5.2.1). We may
suppose, for example, that $s=s_1$ is such that $s_x\ne0$; and since
$\sh{L}_x$ is isomorphic to $\sh{O}_x$, for each $i$ there exists a section
$t_i$ of $\sh{O}_X$ over an open neighborhood $V_i\subset U$ of $x$ such that
$(t_i)_x s_x=(s_i)_x$. Consequently there is an open neighborhood
$V\subset\bigcap_{i=1}^nV_i$ of $x$ such that $t_i s=s_i$ in $V$; in other
words, $\sh{L}|V$ is generated by the \emph{single} section $s|V$. Moreover,
if $z$ is a section over an open subset $W\subset V$ of the kernel of the
homomorphism $\sh{O}_X|V\to\sh{L}|V$ defined by $s$ (5.1.1), then $z_y$
annihilates $\sh{L}_y$ for every $y\in W$, so $z_y=0$ by hypothesis and hence
$z=0$, which completes the proof of our assertion. Finally, consideration of
the tensor product $\sh{L}^{-1}\otimes\sh{L}\otimes\sh{F}$ immediately shows
that $\sh{F}$ is isomorphic to $\sh{L}^{-1}$.

\noindent\textbf{$(\mathbf{Err}_{\mathrm{III}},\,3)$} In
$(0_{\mathrm I},\,7.2.4)$, one must impose the condition that the powers
$\mathfrak J^n$ are \emph{closed} ideals in the admissible ring $A$ in order
for the proposition to be correct. The proof given is incorrect, and the
corrected statement follows from Bourbaki, \emph{Top. gén.}, chap. III, 3rd
ed., \S~3, no.~5, cor. 1 of prop. 9. The $\mathfrak J^n$ must likewise be
assumed \emph{closed} in $A$ in statements $(0_{\mathrm I},\,7.2.5)$ and
$(0_{\mathrm I},\,7.2.6)$.

\noindent\textbf{$(\mathbf{Err}_{\mathrm{III}},\,4)$} After proposition (I,
2.2.5), add: with the notation of (I, 2.2.5), if $\mathfrak J$ is an ideal of
$A$, we denote by $\widetilde{\mathfrak J\sh{F}}$ the
$\sh{O}_X$-submodule $\mathfrak J\widetilde{\sh{F}}$ defined in
$(0,\,4.3.5)$.

\noindent\textbf{$(\mathbf{Err}_{\mathrm{III}},\,5)$} In (I, 3.4.5), on line
1 from the bottom of p. 112, replace ``\emph{a field} $K$'' by ``\emph{an
algebraically closed field} $K$''. On line 1 of p. 113, replace
``\emph{extension}'' by ``\emph{algebraically closed extension}''. On lines 1
to 4 of p. 113, replace the text beginning ``$K$ will be called\ldots{}'' by
the following:

For every point of $X$ with values in a field $K$, $K$ will be called the
\emph{field of values} of the corresponding point; and if $x$ is the locality
of this point, the latter is also said to be \emph{localized at} $x$. This
defines a map $X(K)\to X$ that associates to a point with values in $K$ its
locality.

On lines 7 and 16 of p. 113, delete ``geometric''; on lines 10 and 13 of p.
113, delete ``geometric''. On lines 10 and 11 from the bottom of p. 115,
delete ``\emph{geometric}''.

\noindent\textbf{$(\mathbf{Err}_{\mathrm{III}},\,6)$} At the end of the proof
of (I, 3.6.1), on line 12 from the bottom of p. 117, add: If we set
$X'=X\times_Y\Spec(\sh{O}_y/\mathfrak a_y)$, then, for every point $x\in X'$,
identified by $p$ with a point of $X$, we have
$\sh{O}_{X',x}=\sh{O}_{X,x}/\mathfrak a_y\sh{O}_{X,x}$. Since the question is
local on $X$ and $Y$, we may suppose that $X=\Spec(B)$ and $Y=\Spec(A)$; then
$(B/\mathfrak a_yB)_x=B_x/\mathfrak a_yB_x$ by flatness $(0,\,1.3.2)$.

\noindent\textbf{$(\mathbf{Err}_{\mathrm{III}},\,7)$} In (I, 5.1.1), on line
3 from the bottom of p. 127, replace ``quasi-coherent $\sh{O}_X$-Module'' by
``\emph{quasi-coherent ideal of} $\sh{B}$''.

\noindent\textbf{$(\mathbf{Err}_{\mathrm{III}},\,8)$} After (I, 5.1.10), on
line 16 of p. 131, add:

\emph{Remark (5.1.11).} --- A slight adaptation of the argument used in
(5.1.9) shows that the conclusion remains valid \emph{without assuming a
priori that $X$ is a prescheme}, but

\oldpage[III]{220}
assuming only that $(X,\sh{O}_X)$ is a ringed space in local rings, that
$\sh{J}$ is an ideal of $\sh{O}_X$ such that $\sh{J}^n=0$, that the ringed
space $X_0=(X,\sh{O}_X/\sh{J})$ is an affine scheme, and finally that the
ideals $\sh{J}^k/\sh{J}^{k+1}$ are quasi-coherent $\sh{O}_{X_0}$-Modules. It
suffices to use (1.8.1) (cf. $(\mathbf{Err}_{\mathrm{II}})$) instead of
(2.2.4) in the argument.

\noindent\textbf{$(\mathbf{Err}_{\mathrm{III}},\,9)$} In (I, 5.3.1), on line
11 of p. 132, after ``or $\Delta_X$'', insert ``or $\Delta_\vphi$ if
$\vphi:X\to S$ is the structural morphism''.

\noindent\textbf{$(\mathbf{Err}_{\mathrm{III}},\,10)$} In (I, 5.3.9), the
proof is insufficient, because it does not prove that $\Delta_X(X)$ is
locally closed in $X\times_SX$. To obtain a correct proof, it suffices to use
(4.2.4, $a$)): for every $x\in X$ and every affine neighborhood $U$ of $x$ in
$X$, $U\times_SU$ is an affine neighborhood of $\Delta_X(x)$; taking (5.3.16)
into account (whose proof uses only definition (5.3.1.1)), we are reduced to
proving (5.3.9) when $S=\Spec(B)$ and $X=\Spec(A)$ are affine schemes. It is
then clear from (5.3.1.1) that $\Delta_X$ corresponds to the canonical
homomorphism $A\otimes_BA\to A$ that sends $x\otimes y$ to $xy$; since this
homomorphism is surjective, $\Delta_X$ is then a closed immersion (4.2.3),
which completes the proof.

\noindent\textbf{$(\mathbf{Err}_{\mathrm{III}},\,11)$} In (I, 7.1.15) and
(I, 7.1.16), on lines 7 and 15 from the bottom of p. 158, delete
``\emph{geometric}''.

\noindent\textbf{$(\mathbf{Err}_{\mathrm{III}},\,12)$} Replace (I, 7.4.7)
by: We extend (by abuse of language) the definitions of (7.4.1) to the case
where $X$ is a \emph{reduced} prescheme every point of which admits an open
neighborhood having only a \emph{finite} number of irreducible components. It
then follows from (7.3.4) and (7.4.6) that, for a quasi-coherent
$\sh{O}_X$-Module of finite type $\sh{F}$, to say that $\sh{F}$ is a
\emph{torsion sheaf} is equivalent to saying that $\operatorname{Supp}(\sh{F})$
\emph{contains no irreducible component of $X$}.

\noindent\textbf{$(\mathbf{Err}_{\mathrm{III}},\,13)$} In (I, 10.11.7), on
lines 17 to 23 of p. 205, the end of the proof beginning ``It remains to
prove\ldots{}'' is unnecessary by (10.11.6), since the sheaves in question
are coherent by $(0,\,5.3.5)$.

\noindent\textbf{$(\mathbf{Err}_{\mathrm{III}},\,14)$} In (II, 1.7.8), after
line 10 from the bottom of p. 16, add: when $\sh{E}=\sh{O}_S^n$, we also write
$\mathbf V_S^n$ in place of $\mathbf V(\sh{E})$; if moreover $S=\Spec(A)$ is
affine, we write $\mathbf V_A^n$ in place of $\mathbf V_S^n$; we have
$\mathbf V_A^n=\Spec(A[T_1,\ldots,T_n])$, where the $T_i$ are indeterminates.

\noindent\textbf{$(\mathbf{Err}_{\mathrm{III}},\,15)$} In (II, 1.7.10), on
line 10 of p. 17, delete ``\emph{geometric}''; on lines 12 and 16--17 of p.
17, delete ``geometric''.

\noindent\textbf{$(\mathbf{Err}_{\mathrm{III}},\,16)$} In (II, 1.7.12), add:
We set similarly:
\[
  \mathbf V(\sh{O}_S^n)=\mathbf V_S^n=S[T_1,\ldots,T_n].
\]

\noindent\textbf{$(\mathbf{Err}_{\mathrm{III}},\,17)$} In (II, 4.2.6), on
line 7 of p. 75, delete ``\emph{geometric}''; on line 8 of p. 75, delete
``\emph{geometric}''.

\noindent\textbf{$(\mathbf{Err}_{\mathrm{III}},\,18)$} In (II, 4.6.13), on
line 7 of p. 92, the argument must be made more precise, because with the
notation of (4.4.10) one must prove here that the immersion $\Gamma_f$ is
\emph{quasi-compact}, in order to apply (i bis). By (4.6.4), to prove that
$\sh{L}$ is ample relative to $f$, we may restrict to the case where $Y$ is
affine. Observe also that in each of the hypotheses of (v), $f$ is
quasi-compact (I, 6.6.4). On the other hand, if $g$ is
\oldpage[III]{221}
separated, $\Gamma_f$ is a closed immersion (I, 5.4.3), hence quasi-compact
(I, 6.6.4). If, on the contrary, $X$ is locally Noetherian, then since $Y$ is
affine and $f$ is quasi-compact, the underlying space of $X$ is quasi-compact,
hence Noetherian; and (I, 6.6.4) may again be applied to prove that $\Gamma_f$
is quasi-compact.

\noindent\textbf{$(\mathbf{Err}_{\mathrm{III}},\,19)$} The statement of (II,
5.2.2) is incorrect: conditions $b)$, $c)$, and $c')$ are not local on $Y$
when it is not known whether, for an open subset $U$ of $Y$, a quasi-coherent
$(\sh{O}_X|f^{-1}(U))$-Module is the restriction to $f^{-1}(U)$ of a
quasi-coherent $\sh{O}_X$-Module. Thus, on line 16 of p. 98, replace
``\emph{Let $f:X\to Y$ be a separated quasi-compact morphism}'' by:
``\emph{Let $X$ and $Y$ be two preschemes such that $X$ is a scheme or the
underlying space of $X$ is locally Noetherian, and let $f:X\to Y$ be a
quasi-compact morphism.}'' In the proof, observe that the hypotheses imply
that $f$ is \emph{separated} when $X$ is assumed to be a scheme, by (I, 5.5.5
and 5.5.8); hence (I, 9.2.2) may be applied in both cases.

\noindent\textbf{$(\mathbf{Err}_{\mathrm{III}},\,20)$} In (II, 6.2.3), after
line 18 from the bottom of p. 115, add:

\emph{A morphism $f:X\to Y$ is said to be quasi-finite at a point $x\in X$ if
there exist an affine open neighborhood $V$ of $y=f(x)$ and an affine open
neighborhood $U$ of $x$ such that $f(U)\subset V$ and the morphism $U\to V$
induced by restricting $f$ is quasi-finite. A morphism $f:X\to Y$ is said to
be locally quasi-finite if it is quasi-finite at every point of $X$.}

\noindent\textbf{$(\mathbf{Err}_{\mathrm{III}},\,21)$} In (II, 6.4.3), the
proof is incorrect, since the elements of a finite-type $A$-submodule of $K$
are not necessarily integral over $A$. Replace the last 9 lines of p. 121 by
the following text: Since we may suppose that $E$ is torsion-free, hence
faithful, $E$ is also a faithful $A[u]$-module ($A$ being embedded in the
ring of endomorphisms of $E$). Since $E$ is an $A$-module of finite type, it
follows that $u$ is \emph{integral} over $A$ (Bourbaki, \emph{Alg. comm.},
chap. V, \S~1, no.~1, lemma 1). We immediately conclude that the eigenvalues
of $u\otimes1$ (in an algebraic closure of $K$) are elements \emph{integral}
over $A$, and consequently so are the $\sigma_i(u)$.

\noindent\textbf{$(\mathbf{Err}_{\mathrm{III}},\,22)$} In
$(0_{\mathrm{III}},\,9.1.1)$, on line 18 from the bottom of p. 12, delete
``open''.

\noindent\textbf{$(\mathbf{Err}_{\mathrm{III}},\,23)$} In
$(0_{\mathrm{III}},\,11.7.3)$, on line 10 of p. 43, after $\mathsf C''$, add:
such that for every projective object $P$ of $\mathsf C'$ (resp. every
projective object $P'$ of $\mathsf C''$), the functor
$A'\rightsquigarrow T(P,A')$ (resp. $A\rightsquigarrow T(A,P')$) is exact in
$\mathsf C'$ (resp. $\mathsf C$).

\noindent\textbf{$(\mathbf{Err}_{\mathrm{III}},\,24)$} The statement of
$(0_{\mathrm{III}},\,13.7.7)$ is inaccurate and must be modified as follows:
on line 6 of p. 78, delete ``and
$(\RR^{n+1}T(A_k))_{k\in\mathbf Z}$''; on line 7 of p. 78, replace
``$\RR'^nT(A)$ is an $S$-module of finite type'' by ``$\RR'^nT(A)$ and
$\RR'^{n+1}T(A)$ are $S$-modules of finite type''; on lines 12--13 of p. 78,
delete ``and $p+q=n+1$''. In the proof, on line 23 of p. 78, delete ``and for
$n+1$''.

\noindent\textbf{$(\mathbf{Err}_{\mathrm{III}},\,25)$} In (III, 1.4.15), on
line 6 from the bottom of p. 92, replace ``\emph{of finite type}'' by
``\emph{quasi-compact}''.

\noindent\textbf{$(\mathbf{Err}_{\mathrm{III}},\,26)$} In (III, 2.2.4), on
line 1 of p. 102, replace ``\emph{the supports of $\sh{F}$ and $\sh{H}$ are
proper over $Y$}'' by ``\emph{the support of
$\operatorname{Im}(\sh{F}\to\sh{G})$ is proper over $Y$}''. In the proof,
replace the text of lines 10 to 16, beginning ``This being so\ldots{}'', by:

\oldpage[III]{222}
Set
$\sh{G}_1=\operatorname{Im}(\sh{F}\to\sh{G})
=\operatorname{Ker}(\sh{G}\to\sh{H})$, which is coherent. Since we have the
exact sequence $0\to\sh{G}_1(n)\to\sh{G}(n)\to\sh{H}(n)$ and the functor
$f_*$ is left exact, the sequence
$0\to f_*(\sh{G}_1(n))\to f_*(\sh{G}(n))\to f_*(\sh{H}(n))$ is exact for every
$n$; it therefore suffices to show that, for sufficiently large $n$, the
sequence $f_*(\sh{F}(n))\to f_*(\sh{G}_1(n))\to0$ is exact. In other words, we
may restrict to the case where $\sh{H}=0$ and the support of $\sh{G}$ is
proper over $Y$, hence \emph{closed in $Z$} (II, 5.4.10). Consequently
$\sh{G}'=i_*(\sh{G})$ is a coherent $\sh{O}_Z$-Module such that
$\sh{G}'|X=\sh{G}$. We know from (I, 9.4.3) that there exists a coherent
$\sh{O}_Z$-Module $\sh{F}'$ such that $\sh{F}'|X=\sh{F}$. Moreover, since $X$
is open in $Z$ and $\operatorname{Supp}(\sh{G})\subset X$ is closed in $Z$, it
is immediate that one defines a surjective homomorphism
$u':\sh{F}'\to\sh{G}'$ of sheaves such that $u'|X$ is the given surjective
homomorphism $u:\sh{F}\to\sh{G}$, by taking $u'|U=0$ for every open subset $U$
of $Z$ not meeting $\operatorname{Supp}(\sh{G})$ and $u'|U=u|U$ for every open
subset $U\subset X$. This being so, as at the beginning of the proof of
(2.2.2), we see that we may reduce to the case where $Y$ is affine; it is
therefore enough to show that, for sufficiently large $n$, the homomorphism
$\Gamma(u):\Gamma(X,\sh{F}(n))\to\Gamma(X,\sh{G}(n))$ is surjective. We have
the commutative diagram
\[
\xymatrix@C=35mm@R=12mm{
\Gamma(Z,\sh{F}'(n)) \ar[r]^{\Gamma(u')} \ar[d]_v &
\Gamma(Z,\sh{G}'(n)) \ar[d]^w \\
\Gamma(X,\sh{F}(n)) \ar[r]_{\Gamma(u)} & \Gamma(X,\sh{G}(n))
}
\]
where $v$ and $w$ are the restriction homomorphisms. The definition of
$\sh{G}'$ moreover shows that $w$ is \emph{bijective}; and by (2.2.3),
$\Gamma(u')$ is \emph{surjective} for sufficiently large $n$, so the same is
true of $\Gamma(u)$.

\noindent\textbf{$(\mathbf{Err}_{\mathrm{III}},\,27)$} In (III, 2.2.5),
after line 7 from the bottom of p. 102, add:

(iii) Under the hypotheses of (2.2.4) concerning $X$, $Y$, $f$, and $\sh{L}$,
it is immediate that if $\sh{H}$ is a coherent $\sh{O}_X$-Module whose support
is \emph{proper over $Y$}, an argument analogous to that of (2.2.4) shows
that there exists an integer $N$ such that, for $n\geq N$,
$\RR^if_*(\sh{H}(n))=0$ for every $i>0$. It follows from the exact
cohomology sequence that if $u:\sh{F}\to\sh{G}$ is a homomorphism of coherent
$\sh{O}_X$-Modules such that $\Ker(u)$ and $\Coker(u)$ have supports
\emph{proper over $Y$}, then there exists $N$ such that for $n\geq N$, the
corresponding homomorphism
$\RR^if_*(\sh{F}(n))\to\RR^if_*(\sh{G}(n))$ is \emph{bijective} for every
$i>0$.

\noindent\textbf{$(\mathbf{Err}_{\mathrm{III}},\,28)$} In (III, 4.4.9), on
line 17 from the bottom of p. 137, replace ``\emph{Let $Y$ be a locally
Noetherian integral prescheme}'' by ``\emph{Let $X$ and $Y$ be two locally
Noetherian integral preschemes}''; on line 16 from the bottom of p. 137,
replace ``\emph{of finite type}'' by ``\emph{locally of finite type}''. The
proof is essentially unchanged, since $f^{-1}(y)$ is locally of finite type
over $k(y)$ and is therefore still discrete.

\noindent\textbf{$(\mathbf{Err}_{\mathrm{III}},\,29)$} In (I, 9.3.4), on line
19 from the bottom of p. 173, replace ``\emph{Noetherian}'' by
``\emph{quasi-compact}''; on lines 18 and 19 from the bottom of p. 173,
replace ``\emph{coherent}'' by ``\emph{quasi-coherent of finite type}''. In
the proof, we reduce to the case where

\oldpage[III]{223}
$X=\Spec(A)$ and $\sh{F}=\widetilde M$, where $M$ is an $A$-module of finite
type, and $\sh{J}=\widetilde{\mathfrak J}$, where $\mathfrak J$ is an ideal of
finite type in $A$; the rest of the argument is then unchanged.

\noindent\textbf{$(\mathbf{Err}_{\mathrm{III}},\,30)$} In (I, 9.3.5),
replace lines 1 to 7 from the bottom of p. 173 by the following text:

\emph{Proposition (9.3.5). --- Let $X$ be a prescheme and $\sh{F}$ a
quasi-coherent $\sh{O}_X$-Module of finite type. Then there exist a closed
subprescheme $Y$ of $X$, whose underlying space is equal to
$\operatorname{Supp}(\sh{F})$, and a quasi-coherent $\sh{O}_Y$-Module of
finite type $\sh{G}$ such that, if $j:Y\to X$ is the canonical injection,
$\sh{F}$ is isomorphic to $j_*(\sh{G})$.}

It will suffice to show that the ideal $\sh{J}$ of $\sh{O}_X$,
\emph{annihilator of $\sh{F}$}, is \emph{quasi-coherent}. We then take for $Y$
the closed subprescheme of $X$ defined by $\sh{J}$ (I, 4.1.2), and since
$\sh{J}\sh{F}=0$, $\sh{F}$ is an $(\sh{O}_X/\sh{J})$-Module and we answer the
question by taking $\sh{G}=j^*(\sh{F})$. To see that $\sh{J}$ is
quasi-coherent, we may (the question being local) restrict to the case where
$X=\Spec(A)$ and $\sh{F}=\widetilde M$, where $M$ is an $A$-module generated
by finitely many elements $x_i\ (1\leq i\leq r)$. The ideal $\sh{J}$ is then
the intersection of the annihilators of the $x_i$. But the annihilator of
$x_i$ is the kernel of the homomorphism $\sh{O}_X\to\sh{F}$ corresponding to
the homomorphism $s\mapsto sx_i$ from $A$ to $M$; it is therefore indeed a
quasi-coherent ideal (I, 4.1.1), and every finite intersection of such ideals
is also a quasi-coherent ideal (I, 1.3.10).
